\section{Calibration, implementation error, and the evidence boundary}\label{sec:deploymentcalibration}
The deployment interface is the one stated and proved in
\citet[Section~39]{NoguerAlonso2026Execution}; we record what it requires and how the
present objective enters it. Decision opportunities $k=0,\dots,N$ have states
built from receipt-time information, feasible actions that include an empty
message, conditional mean rewards $r_k$ and next-state kernels $P_k^a$. With a common initial law, observation class, feasible masks and a terminal
reward error at most $\eta_N$, suppose the
fitted model satisfies $\lvert r_k-\widehat r_k\rvert\le\eta_k$ and
$\mathrm{TV}(P_k^a,\widehat P_k^a)\le\epsilon_k$ uniformly, and $S_k$ bounds the
span of fitted continuation values. Then, with $B_N=\eta_N$ and
$B_k=B_{k+1}+\eta_k+\epsilon_kS_{k+1}$ every common feasible policy has
$\lvert J-\widehat J\rvert\le B_0$, and an implemented policy with fitted action
gaps $\zeta_k$ loses at most $2B_0+\widehat{\mathbb E}\sum_k\zeta_k$
\citep[Theorem~39.1]{NoguerAlonso2026Execution}. For $L$ prespecified rows with
$n_\ell$ independent draws each, the simultaneous radii of
\citet[Proposition~39.2]{NoguerAlonso2026Execution} hold with probability $1-\delta$.

Here the reward is the risk-penalized objective of the trajectory certificate or \eqref{eq:allmomentgain}, not cash
P\&L, and the next-state kernel includes the fill fraction of Section~\ref{sec:randomexecution} and the
receipt clock of Section~\ref{sec:productionforecast}. Theorem~\ref{thm:sequentialgain} is a different
statement: it controls realized outcomes given coverage of conditional means,
whereas the table construction supplies such coverage only under independent
row sampling, which dependent book data do not satisfy by themselves. Replay
and off-policy evaluation need the support, invariance and unconfoundedness
conditions of \citet[Section~39.4]{NoguerAlonso2026Execution}; at an unlogged order, two
joint fill/return laws can agree on every observed record and imply opposite
decisions, as the coupling example of Section~\ref{sec:couplingexample} shows.
